\begin{abstract}

Багатосенсорне сприйняття на основі злиття даних~(multi-sensor fusion perception, MSFP) є ключовою технологією для втіленого штучного інтелекту (ШІ), яка обслуговує різноманітні низхідні завдання~(\textit{напр.}, 3D-детектування об'єктів і семантичну сегментацію) та сценарії застосування~(\textit{напр.}, автономне керування та роботи у складі рою).
Останнім часом значні досягнення методів MSFP на основі ШІ було розглянуто у відповідних оглядових працях.
Однак, як ми спостерігаємо після ретельного й детального дослідження, наявні огляди мають низку обмежень.
По-перше, більшість оглядів орієнтовано на окреме завдання або окрему дослідницьку область, наприклад, 3D-детектування об'єктів або автономне керування. Тому дослідникам у суміжних задачах часто складно отримати з них пряму користь.
По-друге, більшість оглядів подають MSFP лише з однієї перспективи мультимодального злиття, не враховуючи різноманітність методів MSFP, як-от багатовидове злиття або злиття часових рядів.
Зважаючи на це, у даній роботі ми прагнемо впорядкувати дослідження MSFP з \emph{не\-залежної від завдання} перспективи, де методи представлено з різних технічних точок зору.
Зокрема, спершу ми вводимо передумови MSFP. Далі розглядаємо методи мультимодального та мульти\-агент\-ного злиття. Далі аналізуємо методи злиття часових рядів. В епоху великих мовних моделей~(LLM) ми також досліджуємо мультимодальні методи злиття на основі LLM (MM-LLM). Нарешті, обговорюємо відкриті виклики та перспективні напрямки для MSFP.
Сподіваємось, що цей огляд допоможе дослідникам зрозуміти важливий поступ у MSFP і дасть можливі відправні точки для майбутніх досліджень.
\end{abstract}

